%============================================================
% CHAPITRE 1 — INTRODUCTION GÉNÉRALE
%============================================================
\chapter{Introduction générale}

\section{Contexte et motivation}

La numérisation des pratiques médicales constitue aujourd'hui un enjeu majeur
pour les systèmes de santé à travers le monde. Au Maroc, ce mouvement de
transformation digitale est particulièrement significatif dans le secteur privé,
où une majorité de cabinets médicaux continuent de fonctionner selon des modes
de gestion traditionnels. Selon les statistiques du Ministère de la Santé marocain,
plus des deux tiers des structures médicales libérales s'appuient encore sur des
supports papier pour la gestion de leurs activités quotidiennes, générant des
inefficacités qui affectent directement la qualité de la prise en charge des patients.

Cette situation se traduit par des problèmes tangibles au quotidien : des dossiers
médicaux difficiles à retrouver ou susceptibles d'être altérés, des plannings mal
maîtrisés entraînant des pertes de temps pour les professionnels et les patients,
une traçabilité financière insuffisante et l'impossibilité de partager rapidement
des informations cliniques entre soignants. Dans ce contexte, la mise en œuvre
de solutions numériques adaptées aux réalités du terrain devient une nécessité
incontournable pour les cabinets qui souhaitent optimiser leur organisation et
élever le niveau de qualité de leur offre de soins.

C'est dans cette perspective que s'inscrit MediCab Pro, une plateforme conçue
en tenant compte des spécificités du marché marocain et construite sur les
fondements des architectures logicielles contemporaines. Ce projet représente
une opportunité unique d'illustrer comment les disciplines DevOps et Cloud
Computing peuvent apporter des réponses concrètes à un domaine aussi exigeant
et réglementé que la santé.

\section{Problématique}

L'analyse approfondie du secteur médical privé marocain révèle un retard numérique
structurel qui se manifeste à travers plusieurs dysfonctionnements récurrents et
mesurables :

\begin{itemize}
  \item \textbf{Fragilité et confidentialité des données médicales} : les documents
    papier sont exposés au risque de détérioration physique, de disparition ou de
    consultation non autorisée, portant atteinte au secret médical et à la protection
    des informations sensibles des patients.

  \item \textbf{Désorganisation opérationnelle} : la gestion manuelle des plannings
    de consultation génère des chevauchements de rendez-vous, des absences non
    signalées et une utilisation sous-optimale du temps médical disponible.

  \item \textbf{Opacité de la situation financière} : sans outil dédié, le suivi
    du chiffre d'affaires, des créances en attente et des remboursements institutionnels
    (CNSS/CNOPS) reste approximatif et chronophage.

  \item \textbf{Manque d'assistance à la décision clinique} : les praticiens ne
    disposent d'aucun outil intelligent pour les seconder dans la formalisation
    des comptes rendus ou l'établissement des ordonnances.

  \item \textbf{Instabilité des déploiements applicatifs} : lorsqu'une solution
    numérique existe, elle est généralement maintenue manuellement, exposée aux
    erreurs et difficile à faire évoluer sans risque de régression.
\end{itemize}

La problématique centrale de ce travail peut donc être formulée ainsi :
\textit{Comment concevoir et mettre en production une plateforme SaaS médicale
sécurisée, scalable et dotée de fonctionnalités intelligentes, s'appuyant sur
une architecture microservices et un pipeline d'intégration et de déploiement
continu, afin de répondre aux besoins opérationnels réels des professionnels
de santé marocains ?}

\section{Objectifs du projet}

MediCab Pro poursuit deux séries d'objectifs complémentaires, fonctionnels
d'une part et techniques de l'autre.

\subsection{Objectifs fonctionnels}

\begin{itemize}
  \item Automatiser et centraliser l'ensemble de la gestion administrative
    d'un cabinet médical : agendas, dossiers patients, consultations et facturation.
  \item Fournir une interface web ergonomique, bilingue (français/arabe) et
    utilisable depuis tout navigateur moderne sans installation préalable.
  \item Intégrer un assistant IA médical (MedAgent) capable d'accompagner
    le médecin en temps réel lors des consultations.
  \item Proposer une architecture multi-tenant permettant l'exploitation simultanée
    par plusieurs cabinets dans un environnement isolé et sécurisé.
  \item Automatiser la production de rapports financiers détaillés pour une
    meilleure pilotage de l'activité du cabinet.
\end{itemize}

\subsection{Objectifs techniques}

\begin{itemize}
  \item Concevoir une architecture microservices robuste et extensible reposant
    sur Spring Boot et Angular 21.
  \item Mettre en place un pipeline CI/CD complet avec Jenkins, Docker et Kubernetes
    pour assurer des livraisons fiables et répétables.
  \item Sécuriser les données médicales sensibles par des mécanismes JWT, OAuth2
    et une gestion des droits basée sur les rôles (RBAC).
  \item Déployer un système d'observabilité complet (Prometheus, Grafana) et
    un déploiement déclaratif GitOps (ArgoCD).
  \item Garantir un taux de disponibilité de 99,9\% grâce à l'auto-scaling
    Kubernetes et aux patterns de résilience (Circuit Breaker, Retry).
\end{itemize}

\section{Méthodologie adoptée}

La conduite du projet MediCab Pro s'est appuyée sur un ensemble de pratiques
méthodologiques reconnues dans l'ingénierie logicielle moderne.

\subsection{Approche Agile}

Le projet a été structuré selon les principes agiles du framework Scrum,
avec des cycles de développement de deux semaines. Cette organisation itérative
favorise une livraison progressive des fonctionnalités, une réactivité aux
retours et une gestion proactive des risques. Chaque itération se conclut
par une démonstration des fonctionnalités développées et une session de
rétrospective visant l'amélioration continue du processus.

\subsection{Domain-Driven Design}

La décomposition du système en microservices s'est appuyée sur les principes
du Domain-Driven Design (DDD), qui préconise d'identifier les contextes
fonctionnels délimités et d'aligner la structure technique sur les domaines
métier réels : gestion de l'identité, suivi des patients, consultations cliniques,
facturation et assistance intelligente.

\subsection{DevOps et automatisation}

L'intégralité du cycle de vie applicatif est prise en charge par un pipeline
CI/CD Jenkins entièrement automatisé. Depuis la soumission du code source jusqu'à
son déploiement sur le cluster Kubernetes, chaque étape est orchestrée : compilation,
exécution des tests, contrôle de la qualité du code, fabrication des images Docker
et synchronisation GitOps via ArgoCD. Cette automatisation réduit significativement
l'exposition aux erreurs humaines et assure la reproductibilité des déploiements.

\subsection{Test-Driven Development}

La maîtrise de la qualité logicielle repose sur une démarche orientée tests :
tests unitaires avec JUnit 5 et Mockito, tests d'intégration avec Testcontainers,
et évaluation de la couverture de code via SonarQube. L'objectif est de maintenir
un taux de couverture supérieur à 80\% pour l'ensemble des microservices critiques.

\section{Structure du rapport}

Le présent document s'articule autour de huit chapitres complémentaires :

\begin{itemize}
  \item \textbf{Chapitre 1 : Introduction générale}\\
    Présentation du cadre général, de la problématique identifiée, des objectifs
    visés, de l'approche méthodologique retenue et de la structure du mémoire.

  \item \textbf{Chapitre 2 : Contexte général du projet}\\
    Élaboration du cahier des charges, spécification des exigences fonctionnelles
    et non fonctionnelles, et description des acteurs du système.

  \item \textbf{Chapitre 3 : Analyse et Conception}\\
    Architecture technique détaillée, modélisation UML (cas d'utilisation,
    séquences, classes) et conception du pipeline DevOps.

  \item \textbf{Chapitre 4 : Réalisation}\\
    Description des outils retenus, des technologies employées et des interfaces
    principales de la solution.

  \item \textbf{Chapitre 5 : Conteneurisation et Déploiement}\\
    Dockerisation des microservices et orchestration sur cluster Kubernetes.

  \item \textbf{Chapitre 6 : Tests et Validation}\\
    Stratégie de test, résultats des tests unitaires, d'intégration et de
    performance.

  \item \textbf{Chapitre 7 : Pipeline CI/CD et Monitoring}\\
    Automatisation du pipeline, supervision avec Prometheus/Grafana et
    déploiement GitOps avec ArgoCD.

  \item \textbf{Chapitre 8 : Conclusion et Perspectives}\\
    Bilan du projet, difficultés rencontrées et axes d'évolution envisagés.
\end{itemize}
